%% USPSC-Introducao.tex

% ----------------------------------------------------------
% Introdução (exemplo de capítulo sem numeração, mas presente no Sumário)
% ----------------------------------------------------------
\chapter[Introdução]{Introdução}
\label{Introdução}


Parte inicial do texto, que contém a delimitação do assunto tratado, objetivos da pesquisa e outros elementos necessários para apresentar o tema do trabalho 




Otimização é uma área da matemática que busca resolver problemas de alta complexidade almejando encontrar a melhor solução que formos capazes para aquele problema, chamada solução ótima. 
A modelagem de problemas, por sua vez, pretende traduzir problemas do mundo real para a linguagem matemática, para conseguirmos futuramente analisá-los e, por exemplo, otimizá-los.

Problemas de corte e empacotamento são problemas otimização que buscam encaixar itens em um dado recipiente, de forma a minimizar desperdício de espaço a ser ocupado. 
Eles podem ser aplicações de problemas da vida real como corte de tecidos na indústria têxtil, corte de itens na automobilística, na indústria de calçados e até na disposição de itens em um contêiner para transporte. 

Nesse trabalho estudamos o problema de corte de itens irregulares em faixa, que consiste em posicionar itens irregulares em uma faixa, sem sobreposição, de forma a minimizar o comprimento usado da faixa. 
Estudamos também ferramentas para resolver esse problema, como NFPs, IFPs, a heurística construtiva \textit{Bottom-Left} e a meta-heurística \textit{BRKGA}.

\textcolor{red}{adicionar referências de trabalhos relacionados}

Na Seção~\ref{sec:irregular} mostramos o problema de corte de itens irregulares em faixa e explicamos principalmente as ferramentas utilizadas para garantir que os itens não se sobreponham e permaneçam no interior da faixa. 
Na Seção~\ref{sec:BL} detalhamos a heurística construtiva \textit{Bottom-Left} implementada para resolvermos o problema apresentado na Seção~\ref{sec:irregular}, juntamente com alguns resultados de seus testes na Seção~\ref{sec:experimentosbl}. 
Na Seção~\ref{sec:BRKGA} detalhamos o funcionamento e implementação da meta-heurística \textit{BRKGA} para a resolução do mesmo problema e analisamos os resultados obtidos na Seção~\ref{sec:experimentosBRKGA}. 
Na Seção~\ref{sec:comp}, comparamos os resultados dos experimentos numéricos de ambas as heurísticas e por fim, na Seção~\ref{sec:cons_finais}, apontamos as considerações finais e possibilidades de trabalhos futuros.